\documentclass[11pt]{article}
\usepackage[a4paper,margin=1.15in]{geometry}
\usepackage{amsmath,amssymb}
\usepackage{enumitem}
\usepackage{array}
\usepackage{xcolor}
\usepackage{mdframed}
\usepackage{hyperref}
\hypersetup{hidelinks}
\title{Pointless Coherence, Part II: Order Without a Centre}
\author{Patrick S.\ Mahon}
\date{October 2026 \quad (draft; reflective, not mathematical)}
\begin{document}
\maketitle

\begin{mdframed}[backgroundcolor=black!4,linewidth=0.4pt]
Part~I asked when local views add up to one world. This part asks where that
question is lived: in ecosystems that hold together with no species on top, and
in societies that hold together with no one in charge. It is a lens, not an
application. Section~5 says plainly where the mathematics reaches and where it
does not yet.
\end{mdframed}

\section{The lizards}
On the hills of central California, male side-blotched lizards come in three
throat colours, each with its own strategy. Orange males hold large territories
and take mates from blue males. Blue males guard their mates closely and so beat
the yellow males. Yellow males sneak into orange territories and beat the
oranges. Orange beats blue, blue beats yellow, yellow beats orange. The
frequencies cycle over the years, and no colour wins (Sinervo and Lively,
\emph{Nature} 1996).

That is Penrose's staircase, alive: every flight goes up, the stairs close into
a loop, and there is no top. Similar three-way cycles are documented among
toxin-producing strains of \emph{E.\ coli}, where the cycle keeps all three
strains alive (Kerr et al., \emph{Nature} 2002), and among encrusting
organisms on coral reefs (Buss and Jackson 1979). Network models suggest that
cycles of this kind can sustain far more species than a ranking could (Allesina
and Levine, \emph{PNAS} 2011).

Ordinary logic says a cycle of ``beats'' is a contradiction, because no ranking
satisfies it. Ecology shows it is not a contradiction at all, but an order. It
is just not a \emph{ranking}. That is the shape of the theory's central finding:
coherent answers with no single world behind them.

\section{Ecology: order from local relations}
\paragraph{Cycles of dominance.} An ecosystem does not consult a global
ranking. Each encounter is local: this lizard, that territory, this season. The
theory's first lesson applies directly. When local relations form a \emph{tree},
they always add up to one global order. When they form a \emph{loop}, they can
be fully coherent with no global order at all. Biodiversity seems to live partly
in those loops.

\paragraph{Steps across scales.} Holling's \emph{panarchy} (Gunderson and
Holling 2002) describes nested cycles at different speeds: a leaf, a stand, a
landscape. Slow scales constrain fast ones, and fast ones occasionally overturn
slow ones. The theory's newest result is that ``beyonds compose''. A triangle can
fail to close even when nothing goes round in a circle, because the step from
first to last must be the composite of the steps in between, and one relation
can ask for a different composite. That rhymes with cross-scale ecology: the
question is not only whether influences go round in a loop, but whether they
\emph{compose} across scales. It is a rhyme, not yet a theorem about ecosystems.

\paragraph{Partial observation.} Ecological knowledge is assembled from plots,
transects, traps and seasons. Each is a local view, correct as far as it goes.
The theory asks of them what a field ecologist should ask of a synthesis: do my
surveys connect in a loop? If not, local agreement is enough. If so, they can
agree pairwise and still describe no single landscape.

\section{Anarchism: order without a sovereign}
\paragraph{The stance is already anarchist.} The theory takes \emph{answers
first}. There is no presupposed ground and no sovereign ``real world'' that the
answers report on. A single world, or a single centre, is something local
answers must \emph{earn}. This mirrors an anarchist account of order: it is not
handed down from a centre, it is made, or not made, from relations.
Feyerabend's ``epistemological anarchism'' (\emph{Against Method}, 1975) makes
the same refusal for science.

\paragraph{Hierarchies are trees.} The forest theorem says that local agreements
arranged as a tree always add up to one global order. A chain of command is a
tree, which is why hierarchies \emph{look} coherent: structurally, they cannot
fail to be. The cost is that coherence is bought by removing every loop, every
relation that does not route through the centre.

\paragraph{Federations are loops.} Kropotkin was a field naturalist before he
was an anarchist. In \emph{Mutual Aid} (1902) he argued from animals, villages,
guilds and medieval cities that cooperation built from local relations is as
real a force as competition. Elinor Ostrom made the case empirically
(\emph{Governing the Commons}, 1990). Swiss alpine pastures, Japanese village
forests, Spanish irrigation huertas and Philippine irrigation societies
governed shared resources for centuries with no state or market in charge. One
of her design principles is \emph{nested enterprises}: overlapping local centres
of rule-making with no single sovereign. Her later \emph{polycentric governance}
generalises it.

In the theory's terms these are webs with loops. Pointless coherence is,
structurally, \emph{order without a sovereign}: every local arrangement fits its
neighbours, and no single centre stands behind the whole.

\paragraph{How commitments chain.} The composition result has a social reading.
Local agreements can each be fine, and fine pairwise, yet still fail to add up.
That happens not because anyone contradicts anyone, but because of how
commitments \emph{chain}: the agreement between the first and last parties
must match the composite of the agreements in between. Federations seem to fail
this way at least as often as through open conflict.

\section{The four questions, lived}
\begin{center}\small
\begin{tabular}{|>{\raggedright\arraybackslash}p{3.9cm}|>{\raggedright\arraybackslash}p{4.6cm}|>{\raggedright\arraybackslash}p{4.6cm}|}
\hline
\textbf{question} & \textbf{in an ecosystem} & \textbf{in a commons or federation}\\\hline
Is there a loop? & cycles of dominance, nutrient cycles, feedbacks & overlapping councils, reciprocal obligations\\\hline
Is there a joint question nobody can ask? & ``who is best?'' when dominance cycles & ``who is in charge?'' in a polycentric order\\\hline
Is an idealisation hiding a contradiction? & equilibrium, carrying capacity, ``in the long run'' & ``the general will'', ``the market'' as one agent\\\hline
Circle, or composition? & do influences compose across scales? & do commitments chain consistently?\\\hline
\end{tabular}
\end{center}

\section{Honest distance}
\begin{center}\small
\begin{tabular}{|>{\raggedright\arraybackslash}p{4.2cm}|>{\raggedright\arraybackslash}p{9.2cm}|}
\hline
\textbf{link} & \textbf{status}\\\hline
lizards, \emph{E.\ coli}, reefs & \emph{literal} as structure: a directed cycle of pairwise dominance is exactly the staircase's shape\\\hline
hierarchy as tree, federation as loop & \emph{structural}: the forest/loop theorem applied to the shape of relations\\\hline
answers-first as anarchist epistemology & \emph{philosophical}: a shared refusal of a presupposed centre\\\hline
panarchy and composition & \emph{a rhyme}: suggestive, not yet a theorem about ecosystems\\\hline
\end{tabular}
\end{center}
The mathematics of Part~I works with idealised infinities. With finitely many
possibilities its own theorem says every coherent pattern \emph{is} a world. The
lizard cycle is coherent as a \emph{dynamics} rather than as a set of answers, and
real commons are finite. So the lived reading needs a different bridge: a
\emph{probabilistic} or \emph{dynamical} version, in which ``locally consistent
but globally unrealisable'' is real for finite systems. Belief propagation on
loopy networks is one such setting. Whether the composition law has a finite,
probabilistic shadow is an open question, and it may be the bridge that matters
most.

\section{What the lived side asks of the mathematics}
\begin{enumerate}[nosep]
  \item \emph{A finite, probabilistic version.} When are local marginals on a
    web of relations consistent yet unrealisable? Does composition, not only
    cycles, govern it?
  \item \emph{Scale as the source of ``beyond''.} In the theory, ``beyond every
    number'' is what creates twist. In ecology and society, the analogue may be
    \emph{separation of scales}. Can that be made precise?
  \item \emph{Dynamics, revisited.} The theory found that motion only moves
    worlds and never creates impossibility. The lizard cycle is a \emph{dynamic}
    order. Is there a version in which a cycle in time, rather than in the
    questions, carries coherence without a top?
  \item \emph{Design.} Ostrom asked which institutions hold together. Could the
    theory ever say which webs of local rules can be coherent without a centre,
    and which must fail by composition?
\end{enumerate}

\bigskip
\noindent\small
Part~I: \texttt{pointless\_coherence\_essence.tex}. The why:
\texttt{WHY.md}. Sources named here are cited from memory of their standard
descriptions. Check them before relying on any detail.
\end{document}
